%%%%%%%%%%%%%%%%%%%%%%%%%%%%%%%%%%%%%%%%%%%%%%%%%%%%%%%%%%%
%
%--------------CAPÍTULO 2 - FUNDAMENTAÇÃO TEÓRICA------------
%
%
\chapter{Fundamentação Teórica}
\label{cap:fundamentacao}

Este capítulo reúne os conceitos usados na proposta descrita no Capítulo~\ref{cap:introducao}. A Seção~\ref{sec:redes-neurais-convolucionais} introduz as redes neurais convolucionais, base de praticamente todos os modelos de visão computacional usados neste trabalho. A Seção~\ref{sec:deteccao-objetos-yolo} define a tarefa de detecção de objetos e descreve a evolução da família YOLO. A Seção~\ref{sec:efficientnet} apresenta a arquitetura EfficientNet, e a Seção~\ref{sec:mistura-especialistas}, o paradigma de Mistura de Especialistas. A Seção~\ref{sec:metricas-avaliacao} formaliza as métricas de avaliação de detectores. A Seção~\ref{sec:validacao-cruzada-vazamento} discute a validação cruzada e o vazamento de dados em conjuntos com dependência temporal. Por fim, a Seção~\ref{sec:explicabilidade-oclusao} apresenta a análise de sensibilidade por oclusão, usada para interpretar as decisões de um dos modelos deste trabalho.

\section{Redes Neurais Convolucionais}
\label{sec:redes-neurais-convolucionais}

As Redes Neurais Convolucionais (\textit{Convolutional Neural Networks}, CNN) são redes neurais artificiais voltadas ao processamento de dados organizados em grade, como imagens \cite{lecun1998gradient}. Em uma rede totalmente conectada, cada neurônio recebe todos os pixels da entrada. Uma CNN, em vez disso, aplica pequenos filtros (\textit{kernels}) que percorrem a imagem e produzem mapas de características (\textit{feature maps}). Essa operação, a convolução, se apoia em duas propriedades da informação visual. A primeira é a localidade: padrões como bordas, texturas e formas costumam depender apenas de uma vizinhança de pixels. A segunda é a invariância à translação: um mesmo padrão deve ser reconhecido em qualquer posição da imagem. Como o mesmo filtro é reutilizado em toda a imagem (compartilhamento de parâmetros), uma CNN tem muito menos parâmetros do que uma rede totalmente conectada equivalente, o que reduz o custo computacional e o risco de sobreajuste.

Uma CNN típica alterna camadas convolucionais com camadas de agrupamento (\textit{pooling}), que reduzem a resolução espacial dos mapas de características e preservam as informações mais relevantes. Ao final, uma ou mais camadas totalmente conectadas combinam as características extraídas em uma predição, como a classe de um objeto. As primeiras camadas tendem a responder a bordas e texturas simples, e as camadas seguintes combinam esses padrões em estruturas cada vez mais abstratas, como partes de objetos e objetos completos. Essa hierarquia é aprendida diretamente dos dados, sem extratores de características projetados à mão, e foi um dos motivos pelos quais as CNNs passaram a dominar a visão computacional depois do sucesso de arquiteturas profundas treinadas em larga escala, como a AlexNet \cite{krizhevsky2012alexnet}. Os detectores da família YOLO (Seção~\ref{sec:deteccao-objetos-yolo}) e a rede EfficientNet (Seção~\ref{sec:efficientnet}), usados neste trabalho, são CNNs organizadas de formas diferentes para resolver tarefas diferentes.

\section{Detecção de Objetos e a Família YOLO}
\label{sec:deteccao-objetos-yolo}

\subsection{A Tarefa de Detecção de Objetos}
\label{subsec:tarefa-deteccao}

Na detecção de objetos, o modelo recebe uma imagem e deve identificar todos os objetos de interesse presentes nela, informando a classe e a localização de cada um. Na classificação de imagens, uma única classe é atribuída à imagem inteira; na detecção, cada objeto é delimitado por uma caixa delimitadora (\textit{bounding box}), em geral um retângulo alinhado aos eixos da imagem e descrito por quatro valores, como as coordenadas do centro, a largura e a altura. O detector também atribui a cada predição um escore de confiança, que indica o quanto o modelo está seguro de que aquela região contém um objeto da classe indicada.

Os detectores baseados em aprendizado profundo costumam ser divididos em duas famílias. Os detectores de duas etapas primeiro geram regiões candidatas a conter um objeto e depois classificam e refinam cada região. Os detectores de uma etapa predizem as caixas e as classes diretamente a partir da imagem completa, em uma única passagem pela rede. A família YOLO, descrita a seguir, segue a segunda abordagem. Detectores de uma etapa tendem a ser mais rápidos e, historicamente, eram um pouco menos precisos que os de duas etapas, diferença que as versões mais recentes reduziram bastante.

\subsection{A Família YOLO}
\label{subsec:familia-yolo}

O YOLO (\textit{You Only Look Once}) \cite{redmon2016yolo} trata a detecção como um único problema de regressão. A imagem de entrada é dividida em uma grade, e cada célula prevê um número fixo de caixas delimitadoras, seus escores de confiança e as probabilidades de classe, tudo em uma única passagem pela rede. Sem uma etapa separada de geração de regiões candidatas, o modelo é treinado e executado de ponta a ponta, o que favorece a velocidade de inferência, um fator importante para o monitoramento em vídeo.

As versões seguintes mantiveram essa formulação e acrescentaram melhorias. O YOLOv2 (ou YOLO9000) incorporou caixas-âncora (\textit{anchor boxes}), formas de referência predefinidas que facilitam a predição de objetos com proporções variadas, e normalização em lote (\textit{batch normalization}), que estabilizou o treinamento \cite{redmon2017yolov2}. O YOLOv3 passou a fazer predições em múltiplas escalas, o que melhorou a detecção de objetos pequenos, e adotou uma rede-base mais profunda, a Darknet-53 \cite{redmon2018yolov3}. O YOLOv4 reuniu um conjunto de técnicas de treinamento e uma rede-base revisada (CSPDarknet), ganhando precisão sem perder velocidade de inferência \cite{bochkovskiy2020yolov4}. As versões mantidas pela Ultralytics abandonaram as caixas-âncora em favor de uma formulação sem âncoras (\textit{anchor-free}), em que o modelo prediz os limites do objeto diretamente a partir de cada ponto da grade, e separaram a predição de classe e a de localização em ramos distintos da rede (cabeça desacoplada). O YOLOv8, usado neste trabalho no teste preliminar com pesos pré-treinados e no treinamento dos detectores próprios (Capítulo~4), segue essa formulação \cite{jocher2023yolov8}.

Os pesos pré-treinados distribuídos com os modelos YOLO costumam ser obtidos no COCO (\textit{Common Objects in Context}) \cite{journals/corr/LinMBHPRDZ14}, uma base de propósito geral com mais de 80 classes de objetos do cotidiano, criada para pesquisas em detecção, segmentação e legendagem de imagens. O COCO não cobre todas as categorias de interesse de aplicações específicas, e caprinos são um desses casos. Por isso, em domínios especializados, os pesos pré-treinados no COCO servem como ponto de partida para o treinamento, e raramente como solução final.

\section{A Arquitetura EfficientNet}
\label{sec:efficientnet}

A EfficientNet é uma família de redes convolucionais de classificação de imagens construída em torno do escalonamento composto (\textit{compound scaling}) \cite{tan2019efficientnet}. Há três formas de aumentar a capacidade de uma rede convolucional: torná-la mais profunda (mais camadas), mais larga (mais filtros por camada) ou aumentar a resolução das imagens de entrada. Trabalhos anteriores costumavam escalonar apenas uma dessas dimensões. A EfficientNet escalona as três em conjunto, de forma equilibrada, controladas por um único coeficiente. O ponto de partida é a EfficientNet-B0, uma arquitetura obtida por busca automática e formada principalmente por blocos convolucionais residuais invertidos com atenção por canal (\textit{squeeze-and-excitation}). O escalonamento da B0 gera as variantes maiores da família (B1 a B7), mais precisas e mais custosas.

A EfficientNet-B0 atinge, no ImageNet, acurácia comparável à de arquiteturas de classificação anteriores de capacidade semelhante, com uma fração dos parâmetros e do custo computacional \cite{tan2019efficientnet}. Isso a torna adequada para tarefas de classificação auxiliares dentro de um sistema maior, em que o classificador não pode pesar no custo total de inferência. É esse o papel da EfficientNet-B0 na arquitetura descrita no Capítulo~4.

\section{Mistura de Especialistas}
\label{sec:mistura-especialistas}

O paradigma de Mistura de Especialistas (\textit{Mixture of Experts}, MoE) divide um problema de aprendizado em subproblemas mais simples, cada um atribuído a um modelo especializado (um \textit{especialista}). Uma rede de roteamento (\textit{gating network}) decide, para cada entrada, qual especialista ou combinação de especialistas deve ser consultada \cite{jacobs1991adaptive}. Dados $K$ especialistas $f_1, \ldots, f_K$ e uma rede de roteamento $g$ que produz um peso $g_k(x)$ para cada especialista a partir da entrada $x$, a formulação original combina as saídas de todos os especialistas por uma soma ponderada:
\begin{equation}
    y(x) = \sum_{k=1}^{K} g_k(x) \, f_k(x).
    \label{eq:moe-soft}
\end{equation}
Nessa formulação, a rede de roteamento e os especialistas são treinados juntos. A rede de roteamento aprende a dividir o espaço de entradas em regiões, e cada especialista se torna mais competente na sua região, sem que todos precisem se ajustar a toda a variabilidade dos dados.

Muitas arquiteturas profundas recentes substituem a combinação ponderada por um roteamento esparso, em que a rede de roteamento seleciona apenas um pequeno subconjunto de especialistas para cada entrada \cite{shazeer2017outrageously}. No limite, apenas um especialista é escolhido (roteamento \textit{top-1}), e a Equação~\ref{eq:moe-soft} dá lugar a
\begin{equation}
    y(x) = f_{k^*(x)}(x), \quad k^*(x) = \arg\max_{k} \, g_k(x).
    \label{eq:moe-hard}
\end{equation}
Como só o especialista selecionado é executado, o custo de inferência cai. Em troca, a decisão de roteamento se torna um ponto único de falha: se a entrada for enviada ao especialista errado, a combinação com os demais não tem como corrigir o erro. A escolha entre roteamento ponderado e esparso, e a granularidade em que o roteamento é feito, dependem do problema; o Capítulo~4 discute essas escolhas para o caso deste trabalho.

\section{Métricas de Avaliação de Detectores de Objetos}
\label{sec:metricas-avaliacao}

Para avaliar um detector, é preciso primeiro decidir quando uma caixa predita corresponde a um objeto anotado (a verdade de referência, ou \textit{ground truth}). Essa correspondência é medida pela Intersecção sobre União (\textit{Intersection over Union}, IoU), a razão entre a área de interseção e a área de união da caixa predita $B_p$ e da caixa de referência $B_g$:
\begin{equation}
    \mathrm{IoU}(B_p, B_g) = \frac{\mbox{Área}(B_p \cap B_g)}{\mbox{Área}(B_p \cup B_g)}.
    \label{eq:iou}
\end{equation}
O IoU vai de 0 (caixas disjuntas) a 1 (caixas idênticas). Fixado um limiar de IoU, tipicamente 0,5, uma predição é um verdadeiro positivo (VP) se houver uma caixa de referência da mesma classe, ainda não associada, com IoU acima do limiar; caso contrário, é um falso positivo (FP). Uma caixa de referência que não é associada a nenhuma predição é um falso negativo (FN) \cite{everingham2010pascalvoc}.

Com essas contagens, definem-se a precisão, fração das predições que estão corretas, e a revocação, fração dos objetos reais que foram detectados:
\begin{equation}
    \mbox{Precisão} = \frac{VP}{VP + FP}, \qquad
    \mbox{Revocação} = \frac{VP}{VP + FN}.
    \label{eq:precision-recall}
\end{equation}
Precisão e revocação costumam variar em sentidos opostos quando se altera o limiar de confiança do detector, por isso é comum resumi-las no F1-\textit{score}, a média harmônica entre as duas:
\begin{equation}
    F1 = 2 \cdot \frac{\mbox{Precisão} \cdot \mbox{Revocação}}{\mbox{Precisão} + \mbox{Revocação}}.
    \label{eq:f1}
\end{equation}

Variando o limiar de confiança, obtém-se uma curva de precisão em função da revocação. A área sob essa curva é a Precisão Média (\textit{Average Precision}, AP) para uma classe e um limiar de IoU fixos. Com várias classes, a Precisão Média Média (\textit{mean Average Precision}, mAP) é a média das AP de cada classe\footnote{No caso de uma classe única, como neste trabalho, a mAP coincide com a AP dessa classe.}. A literatura de detecção costuma reportar a mAP em dois regimes: a mAP50, calculada com um único limiar de IoU de 0,5, e a mAP50-95, média da AP em dez limiares de IoU igualmente espaçados entre 0,5 e 0,95 \cite{journals/corr/LinMBHPRDZ14}. A mAP50-95 é mais exigente, pois penaliza caixas que acertam a presença e a classe do objeto, mas localizam o objeto apenas de forma aproximada.

\section{Validação Cruzada e Vazamento de Dados}
\label{sec:validacao-cruzada-vazamento}

A validação cruzada estima a capacidade de generalização de um modelo, isto é, seu desempenho esperado sobre dados que não foram usados no treinamento. Na variante \textit{K-Fold}, o conjunto de dados é dividido em $K$ subconjuntos (\textit{folds}) de tamanho aproximadamente igual. O modelo é treinado $K$ vezes, e em cada vez um \textit{fold} diferente é reservado para validação e os $K-1$ restantes são usados no treino; o desempenho final resume os $K$ resultados \cite{kohavi1995study}. O procedimento supõe que as amostras sejam independentes entre si. Só assim o desempenho no \textit{fold} de validação diz algo sobre o desempenho em dados novos.

Essa suposição falha quando as amostras têm alguma dependência estrutural, seja temporal, espacial, hierárquica ou de outra natureza, que o particionamento não considera \cite{roberts2017crossvalidation}. Se amostras correlacionadas são distribuídas entre treino e validação, o modelo é avaliado sobre dados que, embora formalmente separados, carregam informação já vista no treino. Esse é o vazamento de dados (\textit{data leakage}) \cite{kaufman2012leakage}. As métricas de validação passam a refletir, em maior ou menor grau, a capacidade do modelo de memorizar amostras parecidas com as de validação, e o resultado típico é um desempenho inflado que não se mantém quando o modelo encontra dados realmente novos. Para reduzir esse risco, o particionamento precisa respeitar a estrutura de dependência dos dados. Uma opção é agrupar as amostras dependentes para que fiquem inteiras de um único lado da partição; outra é impor uma distância mínima entre as amostras de treino e de validação mais próximas. O protocolo adotado neste trabalho, descrito no Capítulo~4, segue a segunda opção.

\section{Explicabilidade por Oclusão}
\label{sec:explicabilidade-oclusao}

Redes neurais profundas costumam ser chamadas de caixas-pretas: produzem predições precisas, mas não é possível entender, a partir dos seus parâmetros, como combinam a entrada para chegar a uma decisão. As técnicas de explicabilidade procuram indicar quais partes da entrada mais influenciaram uma predição específica.

A análise de sensibilidade por oclusão, proposta no contexto da visualização de redes convolucionais \cite{zeiler2014visualizing}, é uma dessas técnicas. Uma pequena região da imagem é substituída por um valor neutro, como um cinza uniforme, e mede-se o quanto a predição do modelo muda em relação à obtida com a imagem original. Repetindo o procedimento para regiões que cobrem toda a imagem, obtém-se um mapa de sensibilidade. As regiões cuja oclusão muda muito a predição são as que mais sustentam a decisão do modelo; as que pouco a alteram são irrelevantes para ela. Ao contrário das técnicas baseadas em gradiente, a oclusão não exige acesso aos parâmetros internos do modelo nem diferenciabilidade em relação à entrada. Basta executar predições, o que permite aplicá-la diretamente a modelos já treinados.

Neste trabalho, a análise foi implementada com a biblioteca Captum \cite{kokhlikyan2020captum}, que oferece oclusão por janelas deslizantes para modelos construídos em PyTorch. O Capítulo~5 detalha sua aplicação ao modelo que identifica o cenário de captura.
